%% IEEE ICRA 2027 -- contributed paper submission (double-anonymous review)
%% ICRA double-column format: IEEEtran conference class, US Letter, 10pt.
%% SUBMISSION VERSION: author names, affiliations and acknowledgements are
%% withheld per the IEEE RAS double-anonymous review rules. Restore them
%% only in the camera-ready version.
\documentclass[conference]{IEEEtran}
\IEEEoverridecommandlockouts

\usepackage{cite}
\usepackage{amsmath,amssymb,amsfonts}
\usepackage{amsthm}
\newtheorem{prop}{Proposition}
\usepackage{graphicx}
\usepackage{booktabs}
\usepackage{textcomp}
\usepackage{xcolor}
\usepackage[colorlinks=true,allcolors=blue]{hyperref}
% Keep author identity out of the PDF document-information dictionary.
\hypersetup{pdfauthor={},pdfsubject={},pdfkeywords={},pdfcreator={}}

% --- convenience macros (match the thesis notation) -------------------------
\newcommand{\R}{\mathbb{R}}
\newcommand{\Lap}{\mathrm{Laplace}}
\newcommand{\loc}{\mu}
\newcommand{\scale}{b}
\newcommand{\KL}{\mathrm{KL}}
\newcommand{\E}{\mathbb{E}}
\DeclareMathOperator*{\argmin}{arg\,min}

\begin{document}

\title{Permutation-Invariant Knowledge Distillation for 
Motion Prediction in Autonomous Vehicles}

%% --- CAMERA-READY AUTHOR BLOCK: do NOT uncomment before acceptance ------
% \author{
% \IEEEauthorblockN{Maria Niki Zografou \quad and \quad Prof.\ Konstantinos Moustakas, \textit{Senior Member, IEEE}}
% \IEEEauthorblockA{\textit{Dept.\ of Electrical and Computer Engineering} \\
% \textit{University of Patras}\\
% Patras, Greece \\
% \{up1096060@ac.upatras.gr, moustakas@ece.upatras.gr\}}}
%% ------------------------------------------------------------------------
\author{
\IEEEauthorblockN{Anonymous Author(s)}
\IEEEauthorblockA{Affiliation withheld for double-anonymous review}}

\maketitle
\thispagestyle{empty}
\pagestyle{empty}

\begin{abstract}
Motion prediction for autonomous vehicles requires models that are both accurate 
and able to run on tight compute restrictions, yet the most successful predictors are too large 
to be deployed and trained on limited hardware. The compression of these models, while retaining accuracy, 
can be achieved through \emph{knowledge distillation}, but modern forecasters 
can emit a \emph{mixture} of trajectories whose modes are ordered arbitrarily. Therefore,
naively distilling mode~$k$ of a student from mode~$k$ of a teacher compares
modes that do not necessarily correspond. We introduce a distillation objective that
scores each teacher mode under the student's \emph{entire} predictive mixture and
therefore needs no mode matching; it is permutation-invariant and remains defined
when teacher and student carry a different number of modes. We show that the
natural mean-target form of this objective (v1) transfers the teacher's mode
\emph{locations} but drives the predicted scales below their calibrated
value, improving geometric accuracy while breaking calibration. A
distribution-matching form (v2) that scores samples carrying the teacher's own
uncertainty scale repairs this. 
On Argoverse 1 with HiVT, distilled students recover roughly $80\%$ 
of the accuracy gap to the next-larger model, matching 
a from-scratch predictor with $3.8\times$ more parameters, while being $15$--$55\times$ 
smaller than the teacher and \emph{preserving calibration}.
\end{abstract}

\begin{IEEEkeywords}
trajectory prediction, knowledge distillation, calibration, mixture density
networks, autonomous driving.
\end{IEEEkeywords}

% ============================================================================
\section{Introduction}
Motion forecasting sits between perception and planning in an autonomous vehicle and
predicts the future positions of every agent in a scene.  That future is inherently
multi-modal; at an intersection an agent may turn or continue straight, so state-of-the-art
models predict several future modes per agent, rather than a single path.

Since driving is a safety-critical task, two properties matter together. 
The forecast must be \emph{accurate}, with its best hypothesis close to the real future,
and \emph{calibrated}, its stated uncertainty must be neither over- nor under-confident.
An accurate but over confident predictor can prove dangerous, if a risk-aware planner 
takes into account the proposed narrow uncertainty intervals, that do not accurately reflect 
the capabilities of the model. Callibration is thus a main safety requirement.

The most accurate of these models are large, while the on-vehicle inference budget
is small; this asymmetry motivates \emph{compression}. Knowledge distillation
(KD)~\cite{hinton2015} can train a small student to imitate a
large teacher, facilitating the creation of small and accurate models. 
It is important to note that applying KD to a multi-modal forecaster cannot be done naively
by matching logits. The model used in this study is HiVT: Hierarchical
vector transformer for multi-agent motion prediction~\cite{hivt}. Following
standard multi-modal forecasting pipelines, HiVT is trained using a 
\emph{winner-takes-all} (WTA) regression loss,
originally proposed in~\cite{mtp2019}, which, on each scene updates only 
the predicted mode closest to the ground truth. WTA produces accurate and diverse modes, 
but it leaves the mapping 
between a mode's \emph{index} and its \emph{meaning} arbitrary: the slot that
encodes ``turn left'' in the teacher may encode ``go straight'' in the student, and
this correspondence varies from scene to scene. A distillation loss that pairs modes by
index therefore regresses one mode onto another. Empirically the identity
pairing is optimal in essentially none of the scenes, so index-aligned KD spends
much of its signal on mismatched pairs.

One could recover a pairing with a Hungarian assignment, at the cost of a hard,
non-differentiable matching step (Sec.~\ref{sec:method}). We take the opposite route
and remove matching entirely.

\noindent\textbf{Contributions.}
\begin{itemize}
\item A \emph{permutation-invariant} distillation objective (without matching) for
Laplace-mixture trajectory predictors that scores every teacher mode under the
student's full mixture density; it is defined for any
teacher/student mode counts.
\item An analysis showing that the natural mean-target form (v1) improves geometry
but \emph{shrinks} the predictive scale and breaks calibration, juxtaposed with a
distribution-matching form (v2) that scores samples carrying the teacher's scale,
and repairs calibration. The v2 objective equals the \emph{forward} Kullback–Leibler 
(KL) divergence from teacher to student, with v1 being its zero-variance limit.
\item Experiments on Argoverse~1 with HiVT backbones: distilled students recover
$\sim 80\%$ of the capacity gap at a fraction of the parameters and memory while
preserving calibration, and inherit the safety-relevant mode-covering behaviour of
the forward KL.
\end{itemize}

% ============================================================================
\section{Related Work}
\subsection{Multi-modal motion forecasting.} Vectorised encoders such as
VectorNet~\cite{vectornet} and transformers such as
HiVT~\cite{hivt} predict a fixed set of $K$ candidate trajectories, typically as a
mixture density~\cite{bishop1994} trained with a WTA or best-of-$K$
objective~\cite{mtp2019,multipath}. We use HiVT as both teacher and student because
its hierarchical, symmetry-aware design remains accurate at small width on a single
GPU.

\subsection{Knowledge distillation.} KD transfers a teacher's soft
targets to a student~\cite{hinton2015} and has been extended to intermediate
features and to structured outputs in object detection. Distillation of
\emph{multi-modal, mixture-valued} predictors is comparatively unexplored. One of
the central obstacles to this approach is that mixture modes have no canonical order.
This paper addresses this challenge, by providing two distinct objectives: 
a mean-target form (v1) and a distribution-matching form (v2). 
Our v2 objective is a distribution-matching loss, 
connecting KD to 
density estimation under the forward Kullback-Leibler divergence (KL).

The closest concept to our approach is Ensemble Distribution Distillation (EnD$^2$)~\cite{endd}, 
which distills an entire distribution over predictions to retain uncertainty. However, our 
setting differs fundamentally: while EnD$^2$ fits a higher-order distribution over 
indexed categorical outputs from an ensemble, we directly match a continuous, 
unordered Laplace mixture from a single  
multi-modal forecaster. Our objective is designed to be entirely matching-free and 
permutation-invariant.

\subsection{Calibration.} A planner consumes the whole predicted
distribution, so a model that is accurate but over-confident is unsafe. We
therefore track calibration (coverage of predicted intervals) alongside geometric
error and examine whether the distillation process preserves it.

% ============================================================================
\section{Preliminaries and Problem}
\label{sec:prelim}
\subsection{Backbone.} HiVT~\cite{hivt} models the scene's interatcions 
in two stages (Fig.~\ref{fig:arch}). A \emph{local} encoder builds, for 
each agent, a symmetry-aware embedding taking into account nearby lane 
vectors and agents and applying temporal attention over its own history. 
All of this is compressed into a single vector per agent. 
A \emph{global} interaction module then exchanges one vector per agent across the
whole scene for long-range context, and a decoder emits the mixture of future trajectories.

\begin{figure}[t]
  \centering
  \includegraphics[width=\linewidth]{figs/hivt_arch_cropped.png}
  \caption{HiVT backbone (reproduced from Zhou~et~al.~\cite{hivt}). Local encoders
  produce one embedding per agent (dashed); a global module exchanges those vectors
  scene-wide (solid) before per-mode decoding. We vary the embedding width $d$ to
  obtain teacher and students.}
  \label{fig:arch}
\end{figure}

\subsection{Mixture output.} For an agent, HiVT predicts a Laplace mixture over the
future trajectory $\mathbf{y}\in\R^{2H}$ ($H$ steps, two coordinates),
\begin{equation}
  \label{eq:mixture}
  p(\mathbf{y}) = \sum_{k=1}^{K} \pi_k\,\Lap\!\left(\mathbf{y}\mid \loc_k,\scale_k\right),
  \quad \pi_k\ge 0,\ \textstyle\sum_{k=1}^{K} \pi_k = 1,
\end{equation}
where mode $k$ has location $\loc_k$ (the trajectory), per-step scale $\scale_k$
(its predicted uncertainty), and weight $\pi_k$. Each component is a product of
$2H$ one-dimensional Laplace densities, $\Lap(y\mid\mu,b)=\tfrac{1}{2b}e^{-|y-\mu|/b}$.

The two levels of the mixture separate two kinds of uncertainty. The weights
$\{\pi_k\}$ and distinct locations $\{\loc_k\}$ capture \emph{which} manoeuvre an
agent takes, while each per-step scale $\scale_k$ captures
the \emph{within-trajectory} spatial uncertainty around a chosen path. A faithful
distillation target must transfer both; a loss that moves the locations but discards
the scales keeps the trajectories while destroying the calibration around the 
predicted path of the trajectory, which is precisely the failure analysed 
in Sec.~\ref{sec:calib}. 
We keep HiVT's Laplace parameterisation unchanged, which corresponds 
to an $L_1$ residual.

\subsection{WTA training.} The regression term updates only the best mode
$k^\star=\argmin_k \sum_t\lVert\loc_{k,t}-y^\star_t\rVert_2$,
\begin{equation}
  \label{eq:reg}
  \mathcal{L}_{\mathrm{reg}} = -\tfrac{1}{H}\textstyle\sum_{t=1}^{H}
    \log\Lap\!\left(y^\star_t\mid \loc_{k^\star,t},\scale_{k^\star,t}\right),
\end{equation}
so the other modes are free to specialise elsewhere. This is what makes the slot
ordering arbitrary. To probe the honesty of the \emph{whole} mixture we use the
mixture negative log-likelihood
\begin{equation}
  \label{eq:mixnll}
  \mathcal{L}_{\mathrm{mix}} = -\log\textstyle\sum_{k=1}^{K}\pi_k\,
    \Lap\!\left(\mathbf{y}^\star\mid\loc_k,\scale_k\right).
\end{equation}


\subsection{The mode-permutation problem.} Because WTA does not set a specific mode order, 
two independently trained models will have different permutations of the same modes. 
As a result, a distillation loss of the form $\sum_k d(\text{student}_k,\text{teacher}_k)$ compares
arbitrary pairs. Since the teacher and student modes are not aligned, 
an index-aligned distillation approach would lead to mismatched comparisons.
A correct objective must be invariant to how either model happens to order its modes.
Fig.~\ref{fig:perm} measures this on the validation set: the per-scene optimal
teacher--student assignment concentrates on a non-identity permutation, and the
identity pairing is optimal in $0\%$ of scenes.

%figure of
\begin{figure}[h]
  \centering
  \includegraphics[width=0.6\linewidth]{figs/mode_perm_optimal.png}
  \caption[Teacher--student mode assignment matrix.]{Per-scene optimal
    teacher--student mode assignment, aggregated over the validation set: cell
    $(i,j)$ counts the scenes whose optimal assignment maps teacher mode $j$ to
    student mode $i$. The mass concentrates on a consistent \emph{non-identity}
    permutation rather than the diagonal, and the full identity pairing is the
    optimal assignment in $0\%$ of scenes (figure title), confirming that the two
    models number their modes in unrelated orders. An index-aligned distillation
    loss would therefore compare mismatched modes on essentially every scene.}
  \label{fig:perm}
\end{figure}


% ============================================================================
\section{Proposed Distillation Framework}
\label{sec:method}
\subsection{A matching-free, permutation-invariant objective}
The student mode $k$ assigns the whole teacher trajectory $\loc^T_f$ the factorised
log-likelihood
\begin{equation}
  \label{eq:traj}
  \ell_{k,f} = \textstyle\sum_{t=1}^{H}\sum_{d=1}^{2}
    \log\Lap\!\left(\loc^T_{f,t,d}\mid \loc^S_{k,t,d},\scale^S_{k,t,d}\right).
\end{equation}
Rather than pair modes, we score $\loc^T_f$ under the student's \emph{entire}
mixture by marginalising over its modes with a numerically stable log-sum-exp,
\begin{equation}
  \label{eq:mixscore}
  \log p_S\!\left(\loc^T_f\right)
   = \log\textstyle\sum_{k=1}^{K_S}\pi^S_k\,
     \Lap\!\left(\loc^T_f\mid\loc^S_k,\scale^S_k\right),
\end{equation}
and average over the teacher's modes weighted by the teacher's own confidence
$\pi^T_f$. This yields the mean-target distillation loss
\begin{equation}
  \label{eq:v1}
  \mathcal{L}_{\mathrm{v1}}
   = -\textstyle\sum_{f=1}^{K_T}\pi^T_f\,
     \log\textstyle\sum_{k=1}^{K_S}\pi^S_k\,
     \Lap\!\left(\loc^T_f\mid\loc^S_k,\scale^S_k\right).
\end{equation}
The log-sum-exp acts as a smooth maximum: teacher mode $f$ is explained by whichever
student mode(s) lie near it, selected by value rather than by index.

\textbf{Property 1 (Permutation invariance).} 
For any objective $\mathcal{L}=-\sum_{f=1}^{K_T}\pi^T_f\,g(p_S,\loc^T_f,\scale^T_f)$ in which the student enters 
\emph{only} through its mixture density $p_S$, reordering the student's or the teacher's modes leaves $\mathcal{L}$
 unchanged. Because the student and teacher index sets enter only through independent sums, no explicit matching is invoked.
Therefore, the objective is well-defined for any $K_S$ and $K_T$, including $K_S \neq K_T$.

\emph{Why not Hungarian matching?} A different approach would be to match 
teacher--student modes in each scene by solving a linear assignment problem and   
then distill the  paired modes. This
reintroduces a hard, non-differentiable step, is undefined when $K_S\neq K_T$, and
collapses ambiguous cases into a single winner-take-all correspondence that can 
potentially flip discontinuously between training steps. Our objective avoids
all three by never forming a pairing.

\subsection{Why the mean target breaks calibration}
\label{sec:calib}
The failure of the mean-target objective (v1) 
is a direct consequence of discarding $\scale^T_f$. By treating the teacher's target 
as a zero-width point mass, the objective removes the penalty for 
under-estimating variance. Scoring a teacher
coordinate $\loc^T_f$ under a student's Laplace density rewards two things: moving the
student's location $\loc^S_k$ toward the target, and shrinking the student's scale
$\scale^S_k$ toward the residual error $|\loc^T_f - \loc^S_k|$. This is visible in the
gradient of the scale, where the only term opposing $\scale\to 0$ vanishes as the 
location term drives the residual to zero:
\begin{equation}
  \label{eq:v1grad}
  \frac{\partial \mathcal{L}_{\mathrm{v1}}}{\partial \scale^S_{k,t,d}}
   = \sum_{f=1}^{K_T}\pi^T_f\, r_{f\to k}
     \!\left(\frac{1}{\scale^S_{k,t,d}}
       - \frac{\lvert \loc^T_{f,t,d}-\loc^S_{k,t,d}\rvert}{(\scale^S_{k,t,d})^2}\right),
\end{equation}
with $r_{f\to k}=\pi^S_k\Lap(\loc^T_f\mid\loc^S_k,\scale^S_k)/p_S(\loc^T_f)$ the
responsibility of student mode $k$ for teacher mean $f$. Setting it to zero, every
stationary scale is
\begin{equation}
  \label{eq:bstar}
  \scale^{S\star}_{k,t,d}
   = \frac{\sum_f \pi^T_f\, r_{f\to k}\,\lvert \loc^T_{f,t,d}-\loc^S_{k,t,d}\rvert}
          {\sum_f \pi^T_f\, r_{f\to k}} .
\end{equation}
The teacher's own within-mode spread $\scale^T_f$ appears \emph{nowhere} in
\eqref{eq:bstar}: v1 matches only the scatter of the teacher's point predictions and
discards the spread each of them carries. The same loss pulls $\loc^S_k$ onto the 
teacher locations it explains, so those residuals shrink toward zero and the
stationary scale $\scale^{S\star}$ shrinks with them, settling well below its
calibrated value. This pathology touches only the scale, leaving the mode weights
$\{\pi^S_k\}$ largely unaffected. The
result is a predictor with better geometric accuracy that reports a lower uncertainty
than justified, becoming systematically over-confident.

\subsection{Distribution matching (v2) and the forward KL}
The fix is to score not the bare teacher location but \emph{samples drawn using the
teacher's own scale}:
\begin{equation}
  \label{eq:v2}
  \mathcal{L}_{\mathrm{v2}} = -\textstyle\sum_{f=1}^{K_T}\pi^T_f\,
    \E_{\mathbf{y}\sim\Lap(\loc^T_f,\scale^T_f)}\!\left[\log p_S(\mathbf{y})\right],
\end{equation}
estimated with $S$ Monte-Carlo draws per teacher mode. The student must now place
density across the whole region the teacher considers plausible. Concretely, repeat
the derivation of \eqref{eq:bstar} for a teacher mode explained by a single student
mode ($r_{f\to k}\!\approx\!1$): the residual is replaced by its expectation under
the teacher's own density, so the stationary scale becomes
\begin{equation}
  \label{eq:v2resid}
  \scale^{S\star}=\E_{y\sim\Lap(\loc^T_f,\scale^T_f)}\lvert y-\loc^S_k\rvert
   = \lvert\Delta\rvert+\scale^T_f\,e^{-\lvert\Delta\rvert/\scale^T_f}
   \;\ge\;\scale^T_f ,
\end{equation}
with $\Delta=\loc^T_{f,t,d}-\loc^S_{k,t,d}$ (indices dropped). The teacher's
spread---absent from \eqref{eq:bstar}---now floors the student's scale, so
$\scale^S\!\to\!0$ is no longer stationary; at a matched location it is recovered
exactly, since a Laplace's mean absolute deviation about its own location is its
scale.

\textbf{Property 2 (Generalization of v1).} 
As $\scale^T_f\to\mathbf{0}$, the Laplace distribution approaches a Dirac delta function
 ($\Lap(\loc^T_f,\scale^T_f)\to\delta_{\loc^T_f}$), and \eqref{eq:v2} naturally reduces to \eqref{eq:v1}.

\textbf{Why Monte Carlo.} Drawing $\mathbf{y}\sim\Lap(\loc^T_f,\scale^T_f)$ is exact
and inexpensive, since the Laplace CDF is invertible in closed form.
But since $p_S$ is a mixture, $\log p_S(\mathbf{y})$ contains a $\log\sum_k(\cdot)$ that
does not decompose across modes, and \eqref{eq:v2} integrates it against the
teacher's own density with no closed form---unlike \eqref{eq:v1}, which evaluates
$\log p_S$ at a single point and needs no integral at all. We therefore use the
sample average $\E_{\mathbf{y}}[\log p_S(\mathbf{y})]\approx\frac{1}{S}
\sum_{s=1}^{S}\log p_S(\mathbf{y}_s)$ over the existing draws (the
estimate adds no sampling cost).

\textbf{Property 3 (Equivalence to the forward KL).} 
Let $p_T(\mathbf{y})=\sum_f\pi^T_f\Lap(\mathbf{y}\mid\loc^T_f,\scale^T_f)$. The v2 objective can be rewritten as:
\begin{equation}
  \label{eq:kl}
  \mathcal{L}_{\mathrm{v2}}
   = -\,\E_{\mathbf{y}\sim p_T}\!\left[\log p_S(\mathbf{y})\right]
   = \KL\!\left(p_T\Vert p_S\right) + H(p_T).
\end{equation}
Since the teacher's entropy $H(p_T)$ is constant with respect to the student, optimizing $\mathcal{L}_{\mathrm{v2}}$ 
is strictly equivalent to minimizing the forward KL divergence $\KL(p_T\Vert p_S)$, 
so v2 is distribution matching in the
strict sense. v1 is the same cross-entropy with the teacher's spread collapsed to
point masses, which is why it transfers locations but discards scales.

\textbf{Why the forward, mode-covering direction.}  We want the student to cover
the teacher's entire probability mass, and therefore all of its plausible modes,
so we use the forward KL $\KL(p_T\Vert p_S)$, which is mode-covering. The reverse
KL is mode-seeking and permits the student to lock onto a single mode. 
In the context of motion prediction, a mode-seeking student would be unsafe: it would ignore the 
possibility of other plausible modes and the planner would not be 
able to prepare for them. The forward direction is also the cheaper one, since it naturally
occurs when sampling from the teacher's distribution, which has been cached to disk --teacher inference does not add any computational cost. 

\textbf{Training objective and implementation.} Both variants enter the student's
loss as a single added term,
\begin{equation}
  \label{eq:total}
  \mathcal{L} = \mathcal{L}_{\mathrm{reg}} + \mathcal{L}_{\mathrm{cls}}
    + \lambda\,\mathcal{L}_{\mathrm{KD}},
  \qquad
  \mathcal{L}_{\mathrm{KD}}\in\{\mathcal{L}_{\mathrm{v1}},\mathcal{L}_{\mathrm{v2}}\},
\end{equation}
with $\mathcal{L}_{\mathrm{reg}}$ the WTA term~\eqref{eq:reg},
$\mathcal{L}_{\mathrm{cls}}$ HiVT's unchanged mode-classification term, and
$\lambda\ge0$ the distillation weight ($\lambda{=}0$ recovers from-scratch
training). $\mathcal{L}_{\mathrm{KD}}$ is averaged over the $2H$ predicted
coordinates, so $\lambda$ sits on the same numerical scale as the regression term.
For v2 we draw $S{=}8$ samples per teacher mode by inverse-CDF reparameterisation.
Both variants compute the student's mixture log-density entirely in log-space with a
numerically stable \texttt{logsumexp} reduction, which prevents the underflow and
\texttt{NaN} gradients that occur when a teacher target lands far from the student's
modes. The frozen teacher's predictions are computed once offline and cached, so
teacher inference never runs during training and distillation adds no architectural
or inference cost to the student. 
% ============================================================================
\section{Experiments}
\label{sec:exp}
\textbf{Setup.} We use Argoverse~1 motion forecasting~\cite{argoverse} ($205{,}942$ train /
$39{,}472$ val scenes; $2$\,s of history, $3$\,s / $H{=}30$ prediction at $10$\,Hz,
$K{=}6$ modes). The teacher is HiVT-128 ($2.56$\,M parameters); students are HiVT-32
and -16 ($170$\,k, $46$\,k). We also compare to the pretrained HiVT-64 model ($653$\,k).
We report geometric best-of-$K$ metrics minADE, minFDE (primary) and miss rate
(MR, endpoint $>2$\,m), and calibration via coverage of the predicted
intervals, the mean scale $b$, and mixNLL. All runs fit on a single consumer GPU. Students are
tuned from scratch before any distillation: HiVT-16 requires $2$ attention heads
(the default $8$ starve each head of dimension at width $16$) and a smaller learning
rate than HiVT-32 ($\eta{=}1\mathrm{e}{-}3$ vs.\ $3\mathrm{e}{-}3$, batch size $128$), 
so the distillation comparison is against a properly tuned rather
than a handicapped baseline. We fix $S{=}8$ Monte-Carlo draws per teacher mode as a
variance/cost trade-off.

Concretely, for each nominal level $p$ we measure the fraction of ground-truth points
falling inside the predicted central $p$-interval; the calibration error is the mean
absolute gap between empirical and nominal coverage over levels, and a curve below the
diagonal (Fig.~\ref{fig:rel}) means the intervals are too narrow.

\textbf{Accuracy degrades but calibration remains healthy.}
Table~\ref{tab:baseline} shows from-scratch baselines: accuracy degrades
faster-than-linearly as width shrinks (the $32\to16$ step costs $+0.45$\,m minFDE
versus $+0.13$\,m for $64\to32$), so the smallest students have the most to recover.
Calibration error is $\le 0.047$ everywhere, and lowest at the smallest widths, so
there is no pre-existing defect for distillation to inherit: any calibration damage
would be introduced by the method itself.

\begin{table}[t]
  \centering
  \caption{From-scratch HiVT baselines on Argoverse~1 validation. Accuracy falls
  faster-than-linearly with width; calibration is already good. The HiVT-16 row is the
  default-seed run; its three-seed mean minFDE is $1.562$ (Table~\ref{tab:seeds16}).}
  \label{tab:baseline}
  \begin{tabular}{lrrrrr}
    \toprule
    Model & Params & minADE$\downarrow$ & minFDE$\downarrow$ & MR$\downarrow$ & calib.\ err \\
    \midrule
    HiVT-128 (teacher) & $2.56$\,M & $0.661$ & $0.969$ & $0.092$ & $0.047$ \\
    HiVT-64            & $653$\,k  & $0.687$ & $1.030$ & $0.103$ & $0.042$ \\
    HiVT-32 (student)  & $170$\,k  & $0.736$ & $1.157$ & $0.122$ & $0.033$ \\
    HiVT-16 (student)  & $46$\,k   & $0.886$ & $1.605$ & $0.199$ & $0.028$ \\
    \bottomrule
  \end{tabular}
\end{table}

\textbf{v1 improves geometry but worsens calibration; v2 repairs it.}
Table~\ref{tab:hivt32} distils into HiVT-32 (single seed per arm). Both v1 and
v2 close most of the geometric gap (minFDE $1.157\to1.050$, $-9.2\%$; minADE
$0.736\to0.698$; MR $0.122\to0.106$). They differ in calibration. Under v1 the
calibration error rises more than five-fold ($0.033\to0.184$) and the mean scale
contracts by $36\%$ ($0.418\to0.268$), so its nominal $90\%$ intervals cover only
$71\%$ of targets. This is the shrinkage predicted by \eqref{eq:bstar}: the
stationary scale is a weighted average of the student's residuals to the
teacher's mean predictions, the location term drives those residuals toward
zero, and the teacher's own spread $\scale^T_f$ never enters. v2 achieves the
same geometry while keeping calibration ($0.033\to0.028$, below both the
baseline and the teacher) and coverage ($0.909$), and it improves mixNLL below
the no-KD baseline. Its mean scale settles at $0.413$, above the teacher's own
$0.358$. This is expected since the student is less accurate, so its residuals are wider
and honest intervals must be wider too. Fig.~\ref{fig:rel} shows the split directly: v1's reliability curve bows below the
diagonal at every nominal level, whereas v2's is indistinguishable from the
from-scratch baseline.

\textbf{Distillation does not collapse the modes.} A separate failure other than
scale shrinkage is \emph{mode collapse}, where a student imitating the teacher
quietly abandons its alternative hypotheses. Two diagnostics, on the v1 sweep 
and the v2 students respectively, show this does not occur. Across a v1
$\lambda$-sweep ($25\%$-data proxy) the mode-weight entropy does not fall,
rising from $1.74$ to $1.76$ nats against the $\ln 6 = 1.792$ ceiling, a $1\%$
move in the direction opposite to collapse, while the mean scale contracts
$32\%$ ($0.542\to0.366$) over the same runs: the pathology of
Sec.~\ref{sec:calib} touches only $\scale$. Geometrically, the mean pairwise
distance between the six predicted mode endpoints is $5.45$\,m for the
v2-distilled HiVT-16 and $5.08$\,m for HiVT-32, bracketing the teacher's own
$5.15$\,m, with fewer than $2\%$ of mode pairs landing within $1$\,m of each
other (below the teacher's $2.6\%$). Distillation tightens the spread relative
to the from-scratch student ($6.37\to5.45$\,m for HiVT-16) but toward the
teacher's more concentrated and more accurate set, not toward zero, which is the
behaviour the mode-covering forward KL of Sec.~\ref{sec:method} predicts.

\begin{table}[t]
  \centering
  \caption{Distillation into HiVT-32 ($\lambda{=}0.5$, single seed per arm). v1 and
  v2 tie on geometry; only v2 preserves calibration.}
  \label{tab:hivt32}
  \begin{tabular}{lrrrr}
    \toprule
     & no-KD & v1 & v2 & teacher \\
    \midrule
    minADE$\downarrow$      & $0.736$ & $0.696$ & $0.698$ & $0.661$ \\
    minFDE$\downarrow$      & $1.157$ & $1.050$ & $1.050$ & $0.969$ \\
    MR$\downarrow$          & $0.122$ & $0.105$ & $0.106$ & $0.092$ \\
    mixNLL$\downarrow$      & $26.6$  & $38.0$  & $24.2$  & $21.1$  \\
    calib.\ err$\downarrow$ & $0.033$ & $0.184$ & $0.028$ & $0.047$ \\
    mean $b$                & $0.418$ & $0.268$ & $0.413$ & $0.358$ \\
    cov@p90                 & $0.903$ & $0.711$ & $0.909$ & $0.894$ \\
    \bottomrule
  \end{tabular}
\end{table}

\begin{figure}[t]
  \centering
  \includegraphics[width=\linewidth]{figs/reliability_v1_v2.png}
  \caption{Reliability (empirical vs.\ nominal coverage) for the three HiVT-32
models of Table~\ref{tab:hivt32} (full data, $\lambda{=}0.5$). A curve below the
diagonal means the predicted intervals are too narrow. v1 is under-covering at
\emph{every} nominal level---its nominal $90\%$ interval captures only $71.1\%$ of
targets---while v2 lies on the from-scratch reference ($0.909$ vs.\ $0.903$).
Distribution matching therefore removes the calibration penalty without giving up
any of v1's geometry.}
\label{fig:rel}
\end{figure}

\textbf{The smallest student benefits most.} Distilling HiVT-16 with v2
($\lambda{=}1.0$) improves minFDE by $21.6\%$ (3-seed mean; $22.7\%$ on the default
single seed), cuts MR by $31\%$, and improves mixNLL by $16\%$, all while
\emph{improving} calibration (Table~\ref{tab:seeds16}). The effect is robust across
seeds: the distilled and non-distilled clouds do not overlap (worst v2 minFDE
$1.24 <$ best no-KD $1.47$), so the gain is not seed noise. The single-run A/B used
the default seed, whose no-KD run sits at the weak end of the baseline spread, so its
$-22.7\%$ headline slightly overstates the $-21.6\%$ seed mean; we quote the seed
mean as the primary figure.

\begin{table}[htbp]
  \centering
  \caption{HiVT-16 multi-seed confirmation (full data, $64$ epochs):
    mean$\pm$sample standard deviation over independent random seeds (no~KD
    $n{=}3$, v2 $n{=}3$). One further from-scratch seed that collapsed into a
    degenerate scale-inflated optimum is excluded as non-converged;
    $\Delta$ is computed on the means. Lower is better
    for all rows. The distillation gain is reproducible and lies well outside the
    seed spread---indeed the two arms do not overlap on any individual seed.}
  \label{tab:seeds16}
  \begin{tabular}{lccr}
    \toprule
    Metric & no KD ($n{=}3$) & v2, $\lambda{=}1.0$ ($n{=}3$) & $\Delta$ \\
    \midrule
    minADE      & $0.875 \pm 0.021$   & $0.771 \pm 0.005$   & $-11.8\%$ \\
    minFDE      & $1.562 \pm 0.080$   & $1.224 \pm 0.018$   & $-21.6\%$ \\
    MR          & $0.189 \pm 0.016$   & $0.130 \pm 0.002$   & $-31.0\%$ \\
    mixNLL      & $34.9 \pm 1.0$      & $29.4 \pm 0.3$      & $-15.7\%$ \\
    calib.\ err & $0.0283 \pm 0.0010$ & $0.0229 \pm 0.0004$ & $-19.1\%$ \\
    \bottomrule
  \end{tabular}
\end{table}

% \begin{table}[t]
%   \centering
%   \caption{v2 distillation into HiVT-16 ($\lambda{=}1.0$), $3$-seed mean$\pm$std per
%   arm, relative to the from-scratch student. Both geometry and calibration improve.
%   Lower is better throughout. ($\Delta$ minFDE also shows the best single run.)}
%   \label{tab:hivt16}
%   \begin{tabular}{lccr}
%     \toprule
%     Metric & no-KD ($n{=}3$) & v2 ($n{=}3$) & $\Delta$ \\
%     \midrule
%     minADE $\downarrow$        & $0.875$  & $0.771$  & $-11.8\%$ \\
%     minFDE $\downarrow$        & $1.562$  & $1.224$  & $-21.6\%$ \,($-22.7\%$ single) \\
%     miss rate $\downarrow$     & $0.189$  & $0.130$  & $-31.0\%$ \\
%     mixNLL $\downarrow$        & $34.9$   & $29.4$   & $-15.7\%$ \\
%     calib.\ err $\downarrow$   & $0.0283$ & $0.0229$ & $-19.1\%$ \\
%     \bottomrule
%   \end{tabular}
% \end{table}

\textbf{The gain concentrates in the hard, far-future tail.} Truncating each
prediction to $1$/$2$/$3$\,s and re-selecting the best mode at each cut-off shows the
v2 minFDE improvement growing monotonically with the horizon: for HiVT-16,
$-7.7\%\to-15.1\%\to-22.7\%$ at $1$/$2$/$3$\,s, and for HiVT-32,
$-3.3\%\to-7.6\%\to-9.2\%$. In the first second the from-scratch student is already
close to the teacher; the teacher's guidance is worth most in the long-horizon,
high-uncertainty part of the trajectory, exactly where limited capacity hurts.

\textbf{Accuracy--capacity frontier.} Fig.~\ref{fig:frontier} places the distilled
students on the parameter--accuracy curve. Distilled HiVT-32 ($1.050$) comes within $2\%$
of the from-scratch HiVT-64 ($1.030$) at $3.8\times$ fewer parameters, recovering
$\sim 84\%$ of the $32\!\to\!64$ gap; HiVT-16 recovers $\sim 83\%$ of the
$16\!\to\!32$ gap, landing within $5.8\%$ of the \emph{un-distilled} HiVT-32 at
$27\%$ of its parameters. Relative to the teacher the students are $15\times$ and
$55\times$ smaller and use $3.6\times$ and $6.3\times$ less GPU memory respectively (Table~\ref{tab:eff}).
That $55\times$ is a parameter and \emph{memory} claim, not a wall-clock one. On CPU the speedup is $3.0\times$ at $bs{=}1$ and
 $5.8$--$11.5\times$ under batching, and on GPU it largely vanishes ($1$--$2.6\times$), 
 since the students never saturate the device and width-independent overhead sets the runtime.

\begin{figure}[t]
  \centering
  \includegraphics[width=\linewidth]{figs/frontier.png}
  \caption{Accuracy--capacity frontier (log parameter axis). Distillation (arrows)
  moves each small student up the curve toward the next larger model at a fraction
  of the parameters.}
  \label{fig:frontier}
\end{figure}

\begin{table}[t]
  \centering
  \caption{Deployment footprint relative to the HiVT-128 teacher. The
  parameter/memory savings are large; latency savings (batched CPU) are real but
  smaller and hardware-dependent.}
  \label{tab:eff}
  \begin{tabular}{lrrr}
    \toprule
    Model & Params & Param$\times$ & GPU mem $\times$ \\
    \midrule
    HiVT-128 (teacher) & $2.56$\,M & $1\times$  & $1\times$ \\
    HiVT-32 (v2)       & $170$\,k  & $15\times$ & up to $\sim$3.6$\times$ \\
    HiVT-16 (v2)       & $46$\,k   & $55\times$ & up to $\sim$6.3$\times$ \\
    \bottomrule
  \end{tabular}
\end{table}

\textbf{Deployment discussion.} The result that matters for deployment is that
calibration survives compression \emph{for free}: v2 adds only a training-time loss
and changes nothing at inference, so the student's compute, memory, and latency are
those of a small model while its predictive honesty matches a large one. This is exactly
what on-vehicle and latency-bound settings demand. The method is agnostic to the
backbone and assumes only a Laplace-mixture head, so it transfers to other mixture
forecasters without change.

\textbf{Ablations and observations.} On the $25\%$-data proxy, v1 geometry
improves steeply from $\lambda{=}0$ and then saturates (minFDE
$1.68\to1.37\to1.39\to1.42$ for $\lambda{=}0/0.25/0.5/1.0$): the proxy ranks
$\lambda\approx0.25$--$0.5$ best, and the difference between them ($1.4\%$) is an
order of magnitude smaller than the $18\%$ gain over $\lambda{=}0$. We carry
$\lambda{=}0.5$ forward for HiVT-32 as a standard value and a
geometry/calibration compromise, and select $\lambda{=}1.0$ for HiVT-16 on
validation minFDE, the smaller student benefiting from a stronger signal. Under
v1, calibration error instead degrades monotonically from the first step
($0.0245\to0.155$ over the same sweep), so $\lambda$ is not directly comparable 
across the two loss variants. Re-running the same sweep under v2 separates the
two effects cleanly: calibration error stays flat in $\lambda$
($0.024/0.022/0.025$ at $\lambda{=}0.25/0.5/1.0$, against $0.0245$ at
$\lambda{=}0$), while geometry keeps improving where v1's has already turned
around (minFDE $1.41/1.30/1.29$ under v2 versus $1.37/1.39/1.42$ under v1).  A
given $\lambda$ therefore means different things under the two objectives, which
is also why the stronger $\lambda{=}1.0$ setting is usable for HiVT-16 under v2:
with the scale-shrinkage pathology removed, extra distillation signal is no
longer paid for in calibration.
% ============================================================================
\section{Conclusion}
We presented a permutation-invariant, matching-free distillation objective for the
multi-modal Laplace-mixture output of a trajectory predictor. Because two
independently trained mixtures number their modes in unrelated orders, any
index-aligned loss compares mismatched modes; scoring the teacher's predictions under
the student's \emph{full} mixture density avoids this, needs no mode assignment, and
stays defined when teacher and student carry different numbers of modes. Of its two
forms, the mean-target one (v1) transfers the teacher's mode locations and weights but
discards its predictive scales, improving best-of-$K$ geometry while degrading
calibration; the distribution-matching one (v2)---the Monte-Carlo forward KL from
teacher to student---restores those scales by scoring samples drawn with
$\scale^T_f$, and on HiVT and Argoverse~1 keeps v1's full geometric gain with
calibration intact. 

Empirically, distilling the HiVT-128 teacher into much smaller students recovers
roughly a full size-class of accuracy at the students' own parameter count: HiVT-32
($6.6\%$ of the teacher's parameters) improves minFDE by $9.2\%$ and lands close to a
from-scratch HiVT-64 with $3.8\times$ more parameters, and HiVT-16 ($1.8\%$) improves
it by about $22\%$ (seed-mean $21.6\%$ over three seeds), approaching an un-distilled
HiVT-32. The gain is larger for the smaller student and concentrated in the
longer-horizon part of the trajectory, and the efficiency dividend is a parameter and
memory saving (up to $6.3\times$ less peak GPU memory) rather than a latency one.

\textbf{Limitations and future work.} Our evaluation is confined to a single teacher,
dataset, and backbone, and uses best-of-$K$ metrics with no closed-loop planning
objective, so the practical value of the retained calibration for a downstream
planner is argued for but not measured. The assignment-based (Hungarian / optimal transport)
alternative is argued against in Sec.~\ref{sec:method} but not benchmarked as a
trained baseline. The HiVT-16 result is supported by a multi-seed study, but the HiVT-32
headline rests on a single converged run, and a matched HiVT-32 seed sweep remains to
be done. Extending the method to multiple backbones and datasets, placing a planner in
the loop, transferring differing mode counts ($K_S \neq K_T$), and modelling epistemic
uncertainty are natural next steps.

% ============================================================================
% AI-use disclosure required by the ICRA 2027 generative-AI policy, which
% covers code as well as text. It names a tool, not a person or a funder, so
% it is compatible with double-anonymous review and STAYS IN the submitted
% version. Personal thanks and funding go here only at camera-ready.
% ============================================================================
\section*{Acknowledgment}
In accordance with the IEEE policy on the use of generative AI, we disclose
that portions of the research code accompanying this work were co-developed
with the assistance of Claude Code (Anthropic). The tool was used as a
programming assistant while implementing and refactoring the distillation
objectives of Sec.~\ref{sec:method} and the evaluation and plotting scripts
behind Sec.~\ref{sec:exp}. It was not used to generate the text of this
paper, its experimental results, or its scientific claims. All
AI-assisted code was reviewed, executed and verified by the authors, who
take full responsibility for the correctness of the implementation and of
the reported results.

\begin{thebibliography}{00}
\bibitem{hinton2015} G.~Hinton, O.~Vinyals, and J.~Dean, ``Distilling the knowledge
in a neural network,'' in \emph{NeurIPS Deep Learning Workshop}, 2015.
\bibitem{hivt} Z.~Zhou, L.~Ye, J.~Wang, K.~Wu, and K.~Lu, ``HiVT: Hierarchical
vector transformer for multi-agent motion prediction,'' in \emph{CVPR}, 2022.
\bibitem{vectornet} J.~Gao \emph{et al.}, ``VectorNet: Encoding HD maps and agent
dynamics from vectorized representation,'' in \emph{CVPR}, 2020.
\bibitem{bishop1994} C.~M.~Bishop, ``Mixture density networks,'' Aston University,
Tech.\ Rep., 1994.
\bibitem{mtp2019} H.~Cui \emph{et al.}, ``Multimodal trajectory predictions for
autonomous driving using deep convolutional networks,'' in \emph{ICRA}, 2019.
\bibitem{multipath} Y.~Chai, B.~Sapp, M.~Bansal, and D.~Anguelov, ``MultiPath:
Multiple probabilistic anchor trajectory hypotheses for behavior prediction,'' in
\emph{CoRL}, 2019.
\bibitem{argoverse} M.-F.~Chang \emph{et al.}, ``Argoverse: 3D tracking and
forecasting with rich maps,'' in \emph{CVPR}, 2019.
\bibitem{endd} A.~Malinin, B.~Mlodozeniec, and M.~Gales, ``Ensemble distribution
distillation,'' in \emph{ICLR}, 2020. arXiv:1905.00076.
\end{thebibliography}

\end{document}